% !TeX root = ../drafts/c-keccak-gkr.tex
\section{Proving Keccak by GKR}\label{annex:keccak}

\subsection{Credits}\label{keccak:credits}

The design is based on Lev Soukhanov's Hashcaster~2~\cite{Hashcaster2}. The packing is a reverse multiplication friendly embedding~\cite{CCXY18}.

\subsection{Overview}\label{keccak:overview}

The prover initially commits to every input state ($1600$ bits) of the $2^{\kbatch}$ permutations being proven. Think of it as $25$ $\K$-valued multilinear polynomials in $\kbatch$ variables. At some point during the protocol, there will be some claim on the output, typically $25$ evaluation claims on the $25$ $\K$-valued polynomials that represent the output state. The goal of GKR~\cite{GKR08} here is to reduce this output claim to an input claim with negligible soundness error.

Keccak-$f$ is the composition of $24$ rounds; each is a linear map followed by chi and the round constant, which we see as a multivariate polynomial of degree two in the state's $1600$ variables over $\F_2$. GKR goes back one round at a time: a \emph{layer} turns a claim on a round's output state into a claim on its input state.

\subsection{Keccak-\texorpdfstring{$f[1600]$}{f[1600]}}\label{keccak:perm}

A state is $1600$ bits: $25$ lanes of $64$ bits, lane $\ell$ in row $\lfloor\ell/5\rfloor$ and column $\ell\bmod5$ of a $5\times5$ grid. Position $j=z+64\ell$ is bit $z$ (\emph{slice} $z$) of lane $\ell$.

For a lane $0\le\ell<25$ and $0\le k<5$, $\along{\ell}{k}$ is the lane in the same row as $\ell$ whose column is that of $\ell$ plus $k$, modulo $5$, and $\along{j}{k}$ is the same slice of that lane:
\[
\along{\ell}{k}=5\lfloor\ell/5\rfloor+\bigl((\ell\bmod5)+k\bigr)\bmod5,\qquad \along{(z+64\ell)}{k}=z+64\,\along{\ell}{k}.
\]
For example $\along{3}{1}=4$ and $\along{4}{1}=0$. The map $j\mapsto\along{j}{k}$ is a bijection with inverse $j\mapsto\along{j}{5-k}$. Chi uses $k=1,2$.

Round $R\in\{1,\dots,24\}$ maps a state $a$ to
\begin{equation}\label{keccak:eq:round}
b=\klin(a),\qquad a'[j]=b[j]+\bigl(1+b[\along{j}{1}]\bigr)\,b[\along{j}{2}]+[\,j<64\,]\,\krc_R[j],
\end{equation}
where $\klin$ is pi $\circ$ rho $\circ$ theta, the product is chi, and $\krc_R$ is FIPS~202's $\mathrm{RC}[R-1]$, added by iota.

\subsection{Packing}\label{keccak:rmfe}

\begin{lemma}[Embedding]\label{lem:rmfe}
There exists $\Ftwo$-linear maps $\kemb:\Ftwo^{16}\to\Ftwo^{\le23}[X]$ and $\kproj:\Ftwo^{\le46}[X]\to\Ftwo^{16}$ with $\kproj\bigl(\kemb(u)\cdot\kemb(v)\bigr)=u\wedge v$ for all $u,v$, $\wedge$ being the coordinatewise product.
\end{lemma}

The proof is in \S\ref{keccak:rmfe-proof}. Let $\kemb_i$ embed the unit vector at $i\in\cube{4}$ and $\kemb_{\allone}=\sum_i\kemb_i$, so that $\kproj(\kemb_{\allone}\cdot\kemb(u))=u$. Each bit of $\kproj(G)$ is a sum of coefficients of $G$, so $\kproj$ applies just as well to a polynomial $G$ with coefficients in $\E$, giving $\kproj(G)\in\E^{16}$. For some weights $\kw\in\E^{16}$, we use the notation $\inner{\kw}{\kproj(G)} = \sum_{i < 16} {\kw}_i{\kproj(G)}_i$ (Definition~\ref{def:inner}).

The key idea is to \emph{pack} the permutations in groups of sixteen. At one position, the bits of a pack's sixteen permutations form a vector $u\in\Ftwo^{16}$, which $\kemb$ turns into a single polynomial. Chi's product of two bits then becomes, for the whole pack, one carry-less product of two such polynomials, from which $\kproj$ reads back the sixteen products.

\subsection{Example: chi on one layer}\label{keccak:idea}

Here we focus only on chi; the linear map and the round constant are left to \S\ref{keccak:layer}. Take $2^{\kbatch}$ permutations in \emph{packs} of sixteen: permutation $\kperm=(i,p)$ is number $i\in\cube{4}$ of pack $p\in\cube{\kbatch-4}$. With $\ain(j;\kperm)$ the bits ($j < 1600$) of permutation $\kperm$ entering chi, those leaving it are
\[
\aout(j;\kperm)=\ain(j;\kperm)+\bigl(1+\ain(\along{j}{1};\kperm)\bigr)\,\ain(\along{j}{2};\kperm).
\]
The verifier holds a claim on a weighted sum of them, with weights $\eq(r,p)$ on the packs ($r\in\E^{\kbatch-4}$), $\kw_i$ within a pack and $\kpw_j$ on the positions:
\[
c=\sum_{p\in\cube{\kbatch-4}}\eq(r,p)\sum_{i\in\cube{4}}\kw_i\sum_{j<1600}\kpw_j\,\aout\bigl(j;(i,p)\bigr).
\]

\paragraph{1. The product polynomial.} In pack $p$ at position $j$, let $S_{p,j}=\kemb\bigl((\ain(j;(i,p)))_{i\in\cube{4}}\bigr)$. Since $\kproj(\kemb_{\allone}\kemb(u))=u$ and $\kproj(\kemb(u)\kemb(v))=u\wedge v$, the polynomial $G_{p,j}=\kemb_{\allone}S_{p,j}+(\kemb_{\allone}+S_{p,\along{j}{1}})\,S_{p,\along{j}{2}}$ projects to the sixteen outputs of pack $p$ at $j$. The prover sends the $47$ coefficients of
\[
G=\sum_{p\in\cube{\kbatch-4}}\eq(r,p)\sum_{j<1600}\kpw_j\,G_{p,j},
\]
and the verifier checks $\inner{\kw}{\kproj(G)}=c$.

\paragraph{2. Random point.} The verifier then evaluates $G$ at a random point $\kpt\in\E$ (soundness error $46/|\E|$) and compares it with the following.
\begin{align*}
G(\kpt)&\qeq\sum_{p\in\cube{\kbatch-4}}\eq(r,p)\sum_{j<1600}\kpw_j\,G_{p,j}(\kpt)\\
&=\sum_{p\in\cube{\kbatch-4}}\eq(r,p)\sum_{j<1600}\kpw_j\Bigl(\kemb_{\allone}(\kpt)\,S_{p,j}(\kpt)+\bigl(\kemb_{\allone}(\kpt)+S_{p,\along{j}{1}}(\kpt)\bigr)\,S_{p,\along{j}{2}}(\kpt)\Bigr)\\
&=\sum_{p\in\cube{\kbatch-4}}\sum_{j<1600}\eq(r,p)\kpw_j\Bigl(\kwsum'\,S_{p,j}(\kpt)+\bigl(\kwsum'+S_{p,\along{j}{1}}(\kpt)\bigr)\,S_{p,\along{j}{2}}(\kpt)\Bigr).
\end{align*}

with $\kwsum'=\kemb_{\allone}(\kpt)$.

\paragraph{3. Sumcheck.} A sumcheck proves the right side above (first on the $\kbatch-4$ bits of $p$, then on the $11$ bits of $j$). It ends at a random pack $\fc^{\mathrm{pack}}\in\E^{\kbatch-4}$ and a random position $\fc^{\mathrm{pos}}\in\E^{11}$, and the verifier has to evaluate:
\[
\eq(r,\fc^{\mathrm{pack}})\;\mle{\kpw}(\fc^{\mathrm{pos}})\;\bigl(\kwsum'\,\kev_0+(\kwsum'+\kev_1)\,\kev_2\bigr),\qquad
\kev_k=\sum_{j<1600}\eq(\fc^{\mathrm{pos}},j)\,S_{\fc^{\mathrm{pack}},\along{j}{k}}(\kpt),
\]
with $S_{\fc,j}=\sum_{p\in\cube{\kbatch-4}}\eq(\fc,p)\,S_{p,j}$. The prover sends $\kev_0,\kev_1,\kev_2$, which lets the verifier conclude the sumcheck. It remains to prove the correctness of $\kev_0,\kev_1,\kev_2$.

\paragraph{4. One claim again.} Each $\kev_k$ is itself a weighted sum of bits. Since $\kemb$ is linear, $S_{p,j}=\sum_i\ain(j;(i,p))\,\kemb_i$, so $S_{p,j}(\kpt)=\sum_i\kw'_i\,\ain(j;(i,p))$ with $\kw'_i=\kemb_i(\kpt)$; and re-indexing the sum over $j$ by the bijection $j\mapsto\along{j}{k}$, of inverse $j\mapsto\along{j}{5-k}$ (\S\ref{keccak:perm}), carries the shift off the bits and onto the weights:
\[
\kev_k=\sum_{p\in\cube{\kbatch-4}}\eq(\fc^{\mathrm{pack}},p)\sum_{i\in\cube{4}}\kw'_i\sum_{j<1600}\eq(\fc^{\mathrm{pos}},\along{j}{5-k})\,\ain\bigl(j;(i,p)\bigr).
\]
The verifier samples $\kbeta$ and batches the three into $c'=\kev_0+\kbeta\,\kev_1+\kbeta^{2}\,\kev_2$, it remains to check:
\[
c'\qeq\sum_{p\in\cube{\kbatch-4}}\eq(\fc^{\mathrm{pack}},p)\sum_{i\in\cube{4}}\kw'_i\sum_{j<1600}\kpw'_j\,\ain\bigl(j;(i,p)\bigr),
\]
with $\kpw'_j=\eq(\fc^{\mathrm{pos}},j)+\kbeta\,\eq(\fc^{\mathrm{pos}},\along{j}{4})+\kbeta^{2}\,\eq(\fc^{\mathrm{pos}},\along{j}{3})$: a claim of the same form we had at the beginning. In other words, after adding the round constants and applying the linear layers, we can reiterate on the next GKR layer.

\subsection{The protocol}\label{keccak:statement}

We prove $2^{\kbatch}$ permutations of Keccak-$f[1600]$ at once, indexed by $\kperm\in\cube{\kbatch}$ low bit first and written $\kperm=(i,p)$ with $i\in\cube{4}$ the low four bits and $p\in\cube{\kbatch-4}$ the pack of \S\ref{keccak:idea}.

The prover commits the inputs as one multilinear $\qkeccak:\cube{\kbatch}\times\cube{5}\to\K$, slot $\ell<25$ of its second coordinate holding input lane $\ell$, bit $z$ being the coefficient of $x^z$; nothing here weighs the seven slots left over, which may hold anything. Its bit slices $Q_z$ satisfy $\qkeccak=\sum_{z<64}Q_z\,x^z$ (Annex~\ref{annex:rs}), so bit $z$ of input lane $\ell$ of permutation $\kperm$ is $Q_z(\kperm,\ell)$.

Let $a_e(j;\kperm)$ be bit $j$ after $e$ rounds \eqref{keccak:eq:round} of permutation $\kperm$, from its committed input, and $\kout{\ell}(\kperm)=\sum_{z<64}a_{24}(z+64\ell;\kperm)\,x^z$ its output lane $\ell$. We are given values $\koutv{\ell}$, claimed to be $\mle{\kout{\ell}}$ at one point $(r_{\mathrm{low}},r_{\mathrm{pack}})\in\E^{4}\times\E^{\kbatch-4}$, and prove
\begin{equation}\label{keccak:eq:statement}
\koutv{\ell}=\mle{\kout{\ell}}(r_{\mathrm{low}},r_{\mathrm{pack}})\quad\text{for all }\ell<25.
\end{equation}

\paragraph{Claims.}\label{keccak:claims} The example's claim, with its weights named: a \emph{claim at depth $d\le24$} is $(c,\kw,r,\kpw)$, with $c\in\E$, \emph{in-pack weights} $\kw\in\E^{16}$, $r\in\E^{\kbatch-4}$ and \emph{position weights} $\kpw\in\E^{1600}$, asserting
\begin{equation}\label{keccak:eq:claim}
c=\sum_{p\in\cube{\kbatch-4}}\ \sum_{i\in\cube{4}}\ \sum_{j<1600}\eq(r,p)\,\kw_i\,\kpw_j\;a_{24-d}\bigl(j;(i,p)\bigr).
\end{equation}
For the first, at depth $0$, the verifier samples $r_{\mathrm{lane}}\in\E^{5}$ once the $\koutv{\ell}$ are in the transcript and sets $\kw_i=\eq(r_{\mathrm{low}},i)$, $r=r_{\mathrm{pack}}$, $\kpw_{z+64\ell}=\eq(r_{\mathrm{lane}},\ell)\,x^{z}$ and $c=\sum_{\ell<25}\eq(r_{\mathrm{lane}},\ell)\,\koutv{\ell}$: by Fact~\ref{fact:eq}, \eqref{keccak:eq:claim} then reads $c=\sum_\ell\eq(r_{\mathrm{lane}},\ell)\,\mle{\kout{\ell}}(r_{\mathrm{low}},r_{\mathrm{pack}})$, a random combination of the claims \eqref{keccak:eq:statement}.

\paragraph{A layer.}\label{keccak:layer} Layer $d<24$ is the example run on round $24-d$: with $b=\klin(a_{23-d})$,
\begin{equation}\label{keccak:eq:chi}
a_{24-d}(j;\kperm)=b(j;\kperm)+\bigl(1+b(\along{j}{1};\kperm)\bigr)\,b(\along{j}{2};\kperm)+[\,j<64\,]\,\krc_{24-d}[j],
\end{equation}
which is chi on the bits of $b$, plus the round constant. Only the two ends change.

\emph{The constant.} It is public and the same in every permutation, and $\eq(r,\cdot)$ sums to one, so its whole share of the claim is $\kwsum\sum_{z<64}\krc_{24-d}[z]\,\kpw_z$ with $\kwsum=\sum_i\kw_i$: the verifier computes that itself and moves it to its own side. Step $1$ is then as in the example, with $S_{p,j}$ and $G$ built on the bits of $b$:
\begin{equation}\label{keccak:eq:projection}
\inner{\kw}{\kproj(G)}+\kwsum\sum_{z<64}\krc_{24-d}[z]\,\kpw_z\ \qeq\ c.
\end{equation}

Steps $2$ and $3$ are unchanged: $\kbatch+7$ sumcheck rounds bind $G(\kpt)$ (\S\ref{sec:prelim}), $\kbatch-4$ over $p$ at two $\E$-elements each, then $11$ over $j$ at three, slice low and lane high ($\kpw_j$ and $S_{p,\along{j}{k}}$ vanish at lanes $25$ to $31$); with $\kw'_i=\kemb_i(\kpt)$, $\kwsum'=\sum_i\kw'_i$ and $\kev_k$ as there, the verifier checks
\begin{equation}\label{keccak:eq:terminal}
\eq(r,\fc^{\mathrm{pack}})\;\mle{\kpw}(\fc^{\mathrm{pos}})\;\bigl(\kwsum'\,\kev_0+(\kwsum'+\kev_1)\,\kev_2\bigr)\ \qeq\ \rc{\kbatch+7}.
\end{equation}

\emph{The linear map.} Step $4$ leaves a claim on the bits of $b$, with weights $\kpw^{\flat}_j=\eq(\fc^{\mathrm{pos}},j)+\kbeta\,\eq(\fc^{\mathrm{pos}},\along{j}{4})+\kbeta^{2}\,\eq(\fc^{\mathrm{pos}},\along{j}{3})$. As $\klin$ is linear, $\sum_j\kpw^{\flat}_j\,b(j;\kperm)=\sum_j(\klin^{\top}\kpw^{\flat})_j\,a_{23-d}(j;\kperm)$: transpose it onto the weights and the claim is on $a_{23-d}$, the layer handing $(c',\kw',\fc^{\mathrm{pack}},\klin^{\top}\kpw^{\flat})$ to depth $d+1$. Nobody ever writes those $1600$ weights out: $\eq(\fc^{\mathrm{pos}},\cdot)$ is a lane table times a slice table and the shifts only move lanes, so a claim's weights stay $\klin^{\top}$ of such a product, which is what $\mle{\kpw}$ is evaluated from.

\paragraph{The opening.}\label{keccak:openings} The claim at depth $24$ is on the committed inputs,
\begin{equation}\label{keccak:eq:in}
c=\sum_{i\in\cube{4}}\kw_i\sum_{o\in\cube{5}}\sum_{z<64}\kpw^{\mathrm{in}}(o,z)\;\mle{Q_z}\bigl(i,r,o\bigr),\qquad
\kpw^{\mathrm{in}}(\ell,\,z)=\kpw_{z+64\ell},
\end{equation}
$\kpw^{\mathrm{in}}$ being zero on slots $25$ and above. A degree-two sumcheck binds $i$ (challenges $\fc^{\mathrm{low}}\in\E^{4}$), then the slot (challenges $\fc^{\mathrm{in}}\in\E^{5}$); the prover sends $s_0,\dots,s_{63}$ and the verifier checks the final claim against $\mle{\kw}(\fc^{\mathrm{low}})\sum_z\mle{\kpw^{\mathrm{in}}}(\fc^{\mathrm{in}},z)\,s_z$, leaving
\begin{equation}\label{keccak:eq:in-claim}
s_z=\mle{Q_z}\bigl(\fc^{\mathrm{low}},r,\fc^{\mathrm{in}}\bigr)\qquad(z<64),
\end{equation}
which ring switching takes (Annex~\ref{annex:rs}).

\subsection{Committing the middle state}\label{keccak:segments}

This is an optimization to get smaller proofs and faster recursion: we also commit to each permutation's state after twelve rounds, its lane $\ell$ in slot $25+\ell$. Fifty lanes a permutation now, so a slot index takes six bits and $\qkeccak$ has $\nkeccak=\kbatch+6$ variables. A permutation becomes two runs of twelve rounds: half $0$ from the committed input to the middle state, half $1$ from there to the output. Writing $a_{m,e}$ for the state after $e$ rounds of half $m<2$ from its committed input, what we prove on top of \eqref{keccak:eq:statement} is that half $0$ lands where half $1$ starts:
\begin{equation}\label{keccak:eq:chain}
a_{0,12}(j;\kperm)=a_{1,0}(j;\kperm)\quad\text{for all }\kperm\text{ and }j<1600.
\end{equation}

Both halves end on a state, and we claim both at the one point $(r_{\mathrm{low}},r_{\mathrm{pack}})$ we are given: half $1$ on the output, where the $\koutv{\ell}$ are, and half $0$ on the middle state, whose value there the prover sends as $c_{\mathrm{st}}$. Reusing that point is the point: one point, one claim. It is random and drawn after $\qkeccak$ is committed, so it pins the middle state as firmly as the output. That is fifty lane claims, which we fold into one by sampling $r_{\mathrm{lane}}$ over the lanes and $r_{\mathrm{seg}}$ over the two halves. The surviving claim carries a coordinate for the half, and twelve layers prove both at once: \eqref{keccak:eq:claim} sums over $m$ too, on $a_{m,12-d}$; layer $d$ crosses round $12(m+1)-d$ of each half, its constant weighted by $\eq(r_{\kbatch-4},m)$; the sumcheck gains a round; and depth $12$ opens as in \S\ref{keccak:openings} over six slot rounds, with $\kpw^{\mathrm{in}}(25m+\ell,z)=\eq(r_{\kbatch-4},m)\,\kpw_{z+64\ell}$: the segment rides the weights, so only $r$'s pack coordinates enter the point.

Only $c_{\mathrm{st}}$ is still the prover's word. Comparing it with slots $25+\ell$ of the commitment, where half $1$ starts, is exactly \eqref{keccak:eq:chain} at a random point: a sumcheck over the slot (challenges $\fc^{\mathrm{out}}$) does it, the prover sending the word $c_{\mathrm{out}}$ and the verifier checking the last running claim against it times the weight that slot carries, which leaves
\begin{equation}\label{keccak:eq:out-claim}
\mle{\qkeccak}\bigl(r_{\mathrm{low}},r_{\mathrm{pack}},\fc^{\mathrm{out}}\bigr)=c_{\mathrm{out}}.
\end{equation}

\subsection{Soundness}\label{keccak:soundness}

Fix $\qkeccak$ before the protocol starts. For every prover, the probability that the verifier accepts with its two final claims \eqref{keccak:eq:out-claim} and \eqref{keccak:eq:in-claim} true of $\qkeccak$, while \eqref{keccak:eq:statement} or \eqref{keccak:eq:chain} fails, is at most
\[
\frac{37\kbatch+902}{|\E|}=\frac{(\kbatch+6)+12\,(3\kbatch+72)+(12+20)}{|\E|},
\]
a term per challenge a false claim must survive (Lemma~\ref{lem:sz}, Fact~\ref{fact:sumcheck}):
\begin{itemize}
\item the start, $\kbatch+6$: $r_{\mathrm{low}}$ and $r_{\mathrm{pack}}$ hide a wrong middle state, $r_{\mathrm{seg}}$ and $r_{\mathrm{lane}}$ a wrong lane among the fifty the first claim combines;
\item a layer, $3\kbatch+72$: $46$ for $\kpt$, as $G$ has degree $46$; $3(\kbatch-3)$ for the pack and segment rounds and $33$ for the $11$ over positions, all of degree three; $2$ for $\kbeta$, in which $c'$ is quadratic;
\item the openings, $12$ and $20$: the six slot rounds of the middle state's opening, then the four instance and six slot rounds of \eqref{keccak:eq:in}, all of degree two.
\end{itemize}
A round message counts as a whole polynomial of bounded degree, of which the compressed messages of \S\ref{sec:prelim} are a special case, and $(r_{\mathrm{low}},r_{\mathrm{pack}})$ is drawn after $\qkeccak$ is committed. The bound holds over any field of characteristic two in which $x^0,\dots,x^{63}$ are independent over $\Ftwo$, as they are in $\E$. Fiat-Shamir and the commitment scheme are not accounted for: the two claims are the output, for ring switching and the PCS to discharge.

\subsection{Proof of the embedding lemma}\label{keccak:rmfe-proof}

\begin{proof}
Take the seven pairwise coprime irreducibles $g$ below, each carrying the bits of $u \in \F_2^{16}$ listed under it, and let $F$ be their product, of degree $23$.
{\setlength{\arraycolsep}{3pt}
\[
\begin{array}{c|c|c|c|c|c|c|c}
g & X & X+1 & X^3+X+1 & X^3+X^2+1 & X^5+X^2+1 & X^5+X^3+1 & X^5+X^3+X^2+X+1\\\hline
\deg g & 1 & 1 & 3 & 3 & 5 & 5 & 5\\
\text{bits} & u_0 & u_1 & u_2,u_3 & u_4,u_5 & u_6,u_7,u_8 & u_9,u_{10},u_{11} & u_{12},u_{13},u_{14}
\end{array}
\]}

Writing $a_0,a_1,a_2$ for the bits $g$ carries, set
\[
  h_g(u)=a_0,\qquad h_g(u)=a_0+(a_0+a_1)X,\qquad h_g(u)=a_0+(a_0+a_1+a_2)X+a_2X^{2}
\]
for $\deg g=1$, $3$ and $5$, so that $h_g(u)$ has degree at most $(\deg g-1)/2$ and takes the values $a_0$ at $0$ and $a_1$ at $1$, with $a_2$ as its coefficient of $X^{2}$. Let $H_u$ be the polynomial of degree below $23$ congruent to $h_g(u)$ modulo every $g$ (Chinese remainder theorem), and put $\kemb(u)=H_u+u_{15}F$, of degree at most $23$. Let $\kproj(P)$ read, $g$ by $g$ in the order of the table, the residue $R=P\bmod g$: its value at $0$ when $\deg g=1$, its values at $0$ and $1$ when $\deg g=3$, those two and its coefficient of $X^{4}$ when $\deg g=5$; the sixteenth bit is the coefficient of $X^{46}$ in $P$. Both maps are $\Ftwo$-linear.

Let $P=\kemb(u)\kemb(v)$ and let $b_0,b_1,b_2$ be the bits of $v$ carried by $g$. Modulo $g$, $P$ is $h_g(u)\,h_g(v)$, of degree at most $\deg g-1$, so that product is the residue $R$ itself. Evaluation being multiplicative, $R(0)=a_0b_0$ and $R(1)=a_1b_1$; and for $\deg g=5$ both factors have degree at most $2$, so the coefficient of $X^{4}$ in $R$ is $a_2b_2$. Those are the first fifteen bits of $u\wedge v$. Finally
\[
  P=H_uH_v+u_{15}FH_v+v_{15}FH_u+u_{15}v_{15}F^{2},
\]
whose first three terms have degree at most $45$ while $F^{2}$ is monic of degree $46$, so the coefficient of $X^{46}$ in $P$ is $u_{15}v_{15}$.
\end{proof}
